\documentclass[12pt]{article}
\usepackage{amsmath, amssymb}
\usepackage{geometry}
\geometry{margin=1in}

\title{Lecture 23: Predicting the Impossible: Tail Bounds and System Reliability}
\author{}
\date{}

\begin{document}

\maketitle

\section*{Context}
This lecture analyzes system reliability under uncertainty using probability tail bounds, motivated by the Amazon Black Friday Sale 2026 scenario \textit{(see pages 2--3)} \cite{lecture}.

\section*{Model Setup}

\subsection*{Definition (Random Variable)}
Let $X$ denote the number of requests (or jobs) arriving in a system within a fixed time interval.

\subsection*{Assumption (Poisson Model)}
From the datacenter model \textit{(page 3)}, arrivals follow a Poisson process:
\[
X \sim \text{Poisson}(\lambda)
\]

\subsection*{Objective}
Compute or bound:
\[
P(X \ge k)
\]
where $k$ is a failure threshold.

\section*{Observation}
From simulation \textit{(page 5)}:
\begin{itemize}
    \item Mean traffic is stable.
    \item Rare tail events ($X \ge k$) cause system failure.
\end{itemize}

\section*{Hierarchy of Bounds}

\subsection*{Insight}
From \textit{page 6}: tighter bounds require more information:
\begin{itemize}
    \item Level 1: Mean $\mu$
    \item Level 2: Variance $\sigma^2$
    \item Level 3: Full distribution (MGF)
\end{itemize}

\section*{Markov's Inequality}

\subsection*{Definition}
For a non-negative random variable $X$:
\[
P(X \ge a) \le \frac{\mathbb{E}[X]}{a}
\]

\subsection*{Requirement}
Only the mean $\mu = \mathbb{E}[X]$.

\subsection*{Property}
\begin{itemize}
    \item Decay rate: $1/k$ (linear)
    \item Loose bound (overestimates probability)
\end{itemize}

\section*{Chebyshev's Inequality}

\subsection*{Definition}
For any random variable $X$:
\[
P(|X - \mu| \ge a) \le \frac{\mathrm{Var}(X)}{a^2}
\]

\subsection*{Equivalent Form}
\[
P(X \ge a) \le \frac{\sigma^2}{(a - \mu)^2}
\]

\subsection*{Requirement}
Mean $\mu$ and variance $\sigma^2$.

\subsection*{Property}
\begin{itemize}
    \item Decay rate: $1/k^2$ (polynomial)
    \item Tighter than Markov
\end{itemize}

\section*{Chernoff Bound}

\subsection*{Definition}
For any $t > 0$:
\[
P(X \ge a) \le \min_{t>0} e^{-ta} M_X(t)
\]
where:
\[
M_X(t) = \mathbb{E}[e^{tX}]
\]

\subsection*{Example Form (Poisson)}
\[
P(X \ge a) \le e^{-\lambda(r\ln r - r + 1)}, \quad r = \frac{a}{\lambda}
\]

\subsection*{Requirement}
Full moment generating function (MGF).

\subsection*{Property}
\begin{itemize}
    \item Decay rate: exponential $e^{-k}$
    \item Tightest bound
\end{itemize}

\section*{Comparison of Decay Rates}

From \textit{page 10}:
\begin{itemize}
    \item Markov: $1/k$ (linear)
    \item Chebyshev: $1/k^2$ (polynomial)
    \item Chernoff: $e^{-k}$ (exponential)
\end{itemize}

\section*{Case Study: Black Friday SLA Breach}

\subsection*{Given (page 11)}
\begin{itemize}
    \item $\lambda = 500$ requests/min
    \item Threshold $k = 600$
\end{itemize}

\subsection*{Objective}
Estimate:
\[
P(X \ge 600)
\]

\section*{Result Comparison (page 12)}

\begin{itemize}
    \item Exact probability: $\approx 1.4\%$
    \item Chernoff bound: $\approx 8.8\%$
    \item Chebyshev bound: $\approx 20\%$
    \item Markov bound: $\approx 90.9\%$
\end{itemize}

\section*{Conclusion}
\begin{itemize}
    \item Markov significantly overestimates risk.
    \item Chebyshev improves accuracy.
    \item Chernoff provides the tightest bound.
\end{itemize}

\section*{Financial Interpretation}

\subsection*{Definition (Waste Multiplier)}
\[
\text{Waste Multiplier} = \frac{\text{Bound Estimate}}{\text{Exact Probability}}
\]

\subsection*{Observation (page 13)}
\begin{itemize}
    \item Markov leads to extreme over-provisioning.
    \item Chernoff minimizes unnecessary cost.
\end{itemize}

\section*{Distribution Sensitivity (page 15)}

\begin{itemize}
    \item For Poisson/Normal: Chernoff performs well.
    \item For heavy-tailed (Gamma): Chernoff may overestimate.
\end{itemize}

\section*{Decision Matrix (page 16)}

\begin{itemize}
    \item Markov:
    \begin{itemize}
        \item Input: Mean
        \item Use: Rough estimates
    \end{itemize}
    \item Chebyshev:
    \begin{itemize}
        \item Input: Mean + Variance
        \item Use: Unknown distribution shape
    \end{itemize}
    \item Chernoff:
    \begin{itemize}
        \item Input: Full MGF
        \item Use: Precise SLA-critical systems
    \end{itemize}
\end{itemize}

\section*{Laplace Distribution Exercise}

\subsection*{Given (page 17)}
\[
X \sim \text{Laplace}(\mu, b)
\]
\[
\sigma = b\sqrt{2}
\]

\subsection*{Objective}
Compute Chernoff bound for:
\[
a = \mu + 3\sigma
\]

\subsection*{Derivation Steps (page 18)}

\begin{itemize}
    \item Step 1:
    \[
    M_X(t) = \mathbb{E}[e^{tX}]
    \]
    \item Step 2:
    \[
    P(X \ge a) \le e^{-ta} M_X(t)
    \]
    \item Step 3:
    Optimize over $t$:
    \[
    \frac{d}{dt} \left(e^{-ta} M_X(t)\right) = 0
    \]
    \item Step 4:
    Substitute optimal $t$ to obtain final bound.
\end{itemize}

\section*{Conclusion}
Tail bounds provide systematic upper bounds on rare event probabilities:
\begin{itemize}
    \item Increasing information yields tighter bounds.
    \item Chernoff bound best approximates true tail behavior.
\end{itemize}

\bibliographystyle{plain}
\begin{thebibliography}{1}
\bibitem{lecture} Lecture Slides: ``Predicting the Impossible: Tail Bounds and System Reliability'' \cite{turn0file0}
\end{thebibliography}

\end{document}